\documentclass[11pt]{article}
\usepackage[margin=1in]{geometry}
\usepackage{amsmath}
\title{An Observation Window at $2048\leq t\leq2049$}
\author{}\date{}
\begin{document}\maketitle
\section{Meaning of the experiment}
The target remains the full nonlinear equation (1.27), with the original synthetic initial condition, coefficients, and periodic box. The autonomous equation allows a local coordinate $s=t-2048$, but the correct local initial value is
\[
F(s,x,y,z)=E(2048+s,x,y,z),\qquad F(0,x,y,z)=E(2048,x,y,z).
\]
Putting the original trigonometric packet at $s=0$ would be a new initial-value problem, not the late evolution of that packet. The current implementation therefore evolves from $t=0$ to $t=2049$ and retains dense field snapshots only on the final interval. It does not skip the preceding evolution.
\section{Implementation and computational cost}
The candidate uses $64^3$ grid points and integrating-factor RK4 with $\Delta t=0.01$, giving 204,900 steps. The linear step factors are cached; the equation and time-stepping formula are unchanged. There are 101 saved complex-field snapshots at spacing $0.01$ on the observation interval. Scalar mass, peak, and spectral-edge diagnostics are recorded every 100 steps throughout the approach and at every saved frame. Progress is printed approximately every 30 seconds.

An atomic checkpoint records the Fourier state, step number, observation frames already collected, diagnostics, and settings approximately every 60 seconds. The same output directory resumes the same calculation; incompatible settings are rejected and a process lock prevents concurrent writers. The default job requests 16 GiB of host memory and one day of execution. A wall-clock budget of 82,000 seconds pauses the calculation before that limit. A paused run must be submitted again with the same directory; a successful scheduler exit does not by itself mean the observation window is complete. If the partition has a shorter limit, reduce both the Slurm time allocation and the internal wall-clock budget accordingly.

Neither a GPU nor a trained PINN is used here. The cluster launch command, from the repository root, is
\begin{verbatim}
git pull
source .venv/bin/activate
sbatch --partition=YOUR_PARTITION cluster/late_window.slurm
\end{verbatim}
The default output directory is
\begin{verbatim}
runs/late-window-2048-2049
\end{verbatim}
Inspect its \texttt{summary.json}. Status \texttt{paused} requests another submission; \texttt{completed\_unvalidated} means the requested fields were saved. Status \texttt{failed\_diagnostics} means a mass-drift or spectral-edge screen exceeded $0.001$ and the run stopped. Nonfinite fields raise an error and retain the last checkpoint. Checkpointing is discrete: an abrupt interruption may require replaying work since the last saved checkpoint.
\section{Validation and limitations}
A reduced-horizon test compares the late-window output with the corresponding slice of an uninterrupted full trajectory. It also pauses before and inside the observation interval and verifies exact agreement with the unpaused calculation. A separate test checks diagnostic failure and misaligned time rejection. These verify implementation, not accuracy at $t=2048$.

The prior $T=8$ experiment failed the sampled peak-refinement screen. This direct late-window run is explicitly exploratory and does not override that failure. Conservation and finite values cannot establish phase accuracy accumulated over thousands of propagation units. Independent grid and time-step refinements must evolve the full preceding trajectory; refining only the final interval cannot remove errors in its starting state. Periodic recurrences also preclude identifying this calculation directly with whole-space decay.

A later local PINN could learn $F(s)$ for $0\leq s\leq1$, using the evolved field as its initial data. The existing trigonometric hard-initial-condition network cannot do this unchanged; its initial-data representation and interface-error checks must be extended first. No late-window PINN training or long-time decay result is claimed here.
\end{document}
